\chapter{Testy i weryfikacja}
Przeprowadzenie weryfikacji potwierdza, że system został wykonany zgodnie z planem i spełnia wszystkie założone wymagania techniczne. Dokumentacja ta stanowi dowód na poprawność działania i niezawodność całej stworzonej architektury.

\section{Testy funkcjonalne}
Prawidłowość realizacji wszystkich procesów logicznych sprawdzono podczas testów funkcjonalnych opisanych w Tabeli \ref{tab:testyfunkcjonalne}, która zawiera zestawienie wymagań i wyników testów.

\begin{table}[ht!]
\centering
\caption{Testy funkcjonalne}
\label{tab:testyfunkcjonalne}
\begin{tabular}{lll}
\toprule
ID wymogu & Opis & Status \\  
\midrule
FR-001      & Serwer uruchomiony i działający w oparciu o SpringBoot.    & Zaliczony \\
FR-002      & Zapis danych do lokalnej bazy SQL na serwerze.    & Zaliczony \\ 
FR-003      & Interfejs użytkownika dostępny przez przeglądarkę i aplikację mobilną.    & Zaliczony \\
FR-004      & Wyświetlanie dzinnika zapisów z datą, godziną i wartościami.    & Zaliczony \\
FR-005      & Integracja z urządzeniami IoT za pomocą RestAPI.    & Zaliczony \\
FR-006      & Implementacja RestAPI do zarządzania danymi i interakcją.    & Zaliczony \\
FR-007      & Przygotowanie dokumentacji w oparcu LaTeX. Wymagania edytorskie.    & Zaliczony \\
FR-008      & Przygotowanie diagramów UML ilustrujących strukturę i zachowanie systemu.    & Zaliczony \\
FR-009      & Dostarczenie działającej aplikacji.    & Zaliczony \\
FR-010      & Utworzenie pliku README.md z opisem projektu i instrukcją instalacji.    & Zaliczony \\
FR-011      & Przygotowanie modelu PCB płytki z mikrokontrolerem ESP32 Wroom.    & Zaliczony \\
\bottomrule 
\end{tabular}
\sourcetab{opracowanie własne}
\end{table}

\section{Testy niefunkcjonalne}
Ocena parametrów jakościowych, takich jak czas odpowiedzi i stabilność, odbyła się w ramach testów niefunkcjonalnych przedstawionych w Tabeli \ref{tab:testyniefunkcjonalne}, potwierdzającej wysoką wydajność oraz poprawną strukturę techniczną architektury.

\begin{table}[ht!]
\centering
\caption{Testy niefunkcjonalne}
\label{tab:testyniefunkcjonalne}
\begin{tabular}{lll}
\toprule
ID wymogu & Opis & Status \\  
\midrule
NFR-001      & Czas odpowiedzi API mniejszy niż 300ms.    & Zaliczony \\
NFR-002      & Dostępność systemu większa niż 99\%.        & Zaliczony \\ 
NFR-003      & Skalowlaność: obsługa min. 1000 jednoczesnych użytkowników.    & Zaliczony \\
NFR-004      & Bezpieczeństwo: szyfrowanie SSL/TLS dla transmisji danych.    & Zaliczony \\
NFR-005      & Kompatybilność: Android 8.0+, przeglądarki wspierające HTML5.	    & Zaliczony \\
NFR-006      & Dostępność zgodna z WCAG 2.1 na poziomie AA.    & Zaliczony \\
\bottomrule 
\end{tabular}
\sourcetab{opracowanie własne}
\end{table}